% kernel/observations-DE.tex

\chapter{Evidenz über den Übergang}
\label{chap:engineering-observations}

Eine Vermessungspolylinie durchquert dasselbe Übergangsintervall. Ihre
Koordinaten ähneln der berechneten Kurve, beantworten aber eine andere Frage.
Der Solver hat abgeleitet, was aus angegebenen Eingaben folgt; die Vermessung
hält Evidenz fest, die an einem physischen Zustand gewonnen wurde.

\begin{definition}[Engineering Observation]
\label{def:engineering-observation}
Eine Engineering Observation ist aufgezeichnete Information über ein
Engineering Object, einen Zustand oder eine Qualifikation zusammen mit den
Bedingungen, unter denen diese Information gewonnen wurde.
\end{definition}

Für die Vermessung umfasst die Provenienz den beobachteten Gegenstand,
Erfassungsmethode und -zeit, Koordinaten- und Realisierungsabhängigkeiten,
Verarbeitungsgeschichte, Unsicherheit und Quellenverantwortung. Beim
historischen Import kennzeichnet Provenienz Quellsystem und Austauschgeschichte.
Beim axtranNew-Ergebnis benennt Operationsprovenienz Eingaben, Algorithmus,
Version und Ableitung. Diese Evidenzwege unterscheiden sich und dürfen nicht in
einer allgemeinen ``Geometriequelle'' verschmelzen.

Mehrere Beobachtungen können voneinander abweichen und dennoch dasselbe
Alignment betreffen. Eine Vermessung wird nicht zur physischen Realität, und
ihre Koordinaten überschreiben nicht den beabsichtigten Zustand. Ein Vergleich
verlangt kompatible Realisierungen und zeitliche Qualifikation. Das entstehende
Residuum ist abgeleitete Evidenz über den angegebenen Vergleich, keine neue
Beobachtung und keine Identitätsentscheidung.

Beobachtung bewahrt, was angetroffen wurde; Provenienz macht seine spätere
Nutzung beurteilbar. Die nächste Frage lautet, auf welche Regeln, angenommenen
Ableitungen und früheren Entscheidungen bei der Bewertung des Kandidaten
berechtigterweise vertraut werden darf.
